% Research proof for the explicit product scope only.
\begin{proof}
The argument constructs the divisor functional from actual
compactly supported test forms. It does not restrict \(T\) to a
curve, take a current slice, or pull \(T\) through a field extension.

\emph{The semistable neighbourhood and its individual branches.}
The formal product
\[
 \operatorname{Spf}R\langle x,y\rangle/(xy-\pi)
       \times_R\widehat{\mathcal C}
\]
is strictly semistable of relative dimension two. Its two
special-fibre components meet along the proper smooth curve
\(C_0\). The space \(X\) is the ordinary open part of its generic
fibre with \(0<v(x)<\ell\), using \(z=x\). This gives the
advertised codimension-one skeletal face with proper residual
curve. The product structure and properness of \(C_0\) are
explicit hypotheses, not conclusions from a general etale chart.

For a closed point \(p\), perfection of \(k\) makes
\(\kappa(p)/k\) finite separable. The unramified extension \(R_p/R\)
exists by the henselian lifting of finite separable residue
extensions and has degree \(d_p=[\kappa(p):k]\).
Hensel lifting embeds \(R_p\) into
\(\widehat{\mathcal O}_{\mathcal C,p}\), with the specified residue
embedding. Choose \(w_p\) lifting a uniformizer of the regular
local ring \(\mathcal O_{C_0,p}\). Smoothness implies
\[
 \widehat{\mathcal O}_{\mathcal C,p}\simeq R_p[[w_p]].
\]
Indeed, \(\pi,w_p\) generate the maximal ideal, the residue map is
an isomorphism, and successive approximation gives a surjection
from \(R_p[[w_p]]\) onto the completed local ring. Both rings are
regular local domains of dimension two, so its kernel is zero.

The generic fibre of this completion is
\(\{|w_p|<1\}\) over \(K_p\): a formal power series with bounded
coefficients converges on every closed disc of radius less than
one, and these discs exhaust the tube. This also follows directly
by using the affinoid subdomains
\(\{|w_p|^N\leq|\pi|\}\), \(N\geq1\).
Thus an outer collar is the actual \(K\)-analytic space
\[
 V_p=A_{K_p}\times_{K_p}
        \{e^{-\delta_p}<|w_p|<1\},
\]
viewed over \(K\). Distinct \(p\)'s give disjoint ordinary open
source domains. No decomposition into \(d_p\) ordinary source
sheets is asserted.

The tropical map
\[
 q_p:V_p\longrightarrow Q_p,\qquad
        q_p=(v(z),v(w_p)),
\]
has compact inverse image of every compact subset of \(Q_p\).
Such an inverse image is closed in a closed polyannulus over
the complete field \(K_p\), and that polyannulus is compact.
Consequently all pullbacks of compactly supported target test
forms have compact analytic support in \(V_p\). Their extensions
by zero to \(X\) are legitimate smooth test forms. This is actual
properness on compact target sets, not an appeal to a generic
finite map or to ordinary-open covering sheets.
Compare the exact pushforward requirement in
\cite[Section 7.1.11, p.~159]{CLD25}.

\emph{A finite constant on each collar.}
Let
\[
 \Omega=\mathrm d'u\wedge\mathrm d''u\wedge
          \mathrm d's\wedge\mathrm d''s
\]
on \(Q_p\). Define a distribution by
\[
 \nu_p(h)=\langle T,q_p^*(h\Omega)\rangle.
\]
It is positive. For tests supported in a fixed compact set,
choose a nonnegative smooth cutoff equal to one there.
Positivity bounds \(|\nu_p(h)|\) by the sup norm of \(h\) times
the finite value on this cutoff. Hence \(\nu_p\) is a positive
Radon measure, finite on compact sets. This is the ordinary
Euclidean positive-distribution argument, consistent with
\cite[Sections 7.2.6 and 7.2.8--7.2.9, pp.~161--162]{CLD25}.

Deleting \(\mathrm d'u\), respectively \(\mathrm d's\), from
\(\Omega\), with the appropriate sign, and using
\(\mathrm d'T=0\), gives
\[
 \nu_p(\partial_u h)=\nu_p(\partial_s h)=0
 \qquad(h\in C_c^\infty(Q_p)).
\]
All pulled-back primitive forms have compact support by the
preceding paragraph. On a ball compactly contained in \(Q_p\),
convolve \(\nu_p\) with a Euclidean mollifier. Both derivatives
of the convolution are zero, so it is constant. Testing with
one function of integral one shows that the constants have a
finite limit as the mollifier radius tends to zero. Thus
\(\nu_p=b\,\mathrm du\,\mathrm ds\) on that ball. Overlapping
balls give the same \(b\); connectedness of \(Q_p\) gives a
single finite \(b_T(p)\geq0\). This proves
\eqref{eq:p213-residual-branch-coefficient}. Only the target
rectangle was mollified.

Restriction to a smaller outer collar does not change this
constant. If \(w'_p\) is another formal parameter, write it in
\(R_p[[w_p]]\) as
\[
 w'_p=a_0+a_1w_p+a_2w_p^2+\cdots,\qquad
 a_0\in\pi R_p,\quad a_1\in R_p^\times,\quad a_j\in R_p.
\]
For \(|w_p|\) sufficiently close to one from below, the linear
term strictly dominates the constant and higher terms.
Therefore \(|w'_p|=|w_p|\) on an entire ordinary outer collar,
not just at its Gauss points. The two tropical top forms there
are equal as actual forms, so the two constants coincide.
Choices of the coefficient-algebra lift are compared through the
specified residue identification; henselian uniqueness gives the
same unramified coefficient algebra. This proves the stated
independence.

\emph{Lifting an arbitrary residual rational function.}
Fix \(f\in k(C_0)^\times\). The local ring of \(\mathcal C\) at
the generic point of \(C_0\) is a DVR with residue field
\(k(C_0)\). Lift \(f\) to a unit \(F\) of that ring. It is a
rational function on \(C\), with vertical valuation zero, and
reduces to \(f\). No globally invertible lift on all of
\(C^{\mathrm{an}}\) is asserted.

At \(p\), express \(F=a/b\) in
\(\operatorname{Frac}(\mathcal O_{\mathcal C,p})\).
After cancelling equal powers of \(\pi\), both reductions
\(\bar a,\bar b\) are nonzero. In \(R_p[[w_p]]\), let \(i\) and
\(j\) be the orders of their reductions. Their coefficients below
these orders lie in \(\pi R_p\), the coefficients at these
orders are units, and the higher coefficients have absolute
value at most one. For radii sufficiently close to one, the
terms of orders \(i\) and \(j\) strictly dominate all other
terms. On a whole outer collar, therefore,
\begin{equation}
\label{eq:p213-residual-lift-collar}
 |F|=|w_p|^{i-j},
 \qquad
 t:=v(F)=m_ps_p,\qquad
 m_p=i-j=\operatorname{ord}_p(f).
\end{equation}
This proves the exact slope on the ordinary collar, including
the assertion that \(t\) is identically zero there when \(m_p=0\).

Let \(B\subset C_0\) be the finite set of specializations of the
zeros and poles of \(F\) on \(C\). Finiteness follows because
their closures are finitely many horizontal divisors in the
proper regular model and meet the special fibre in finitely
many points. Since the vertical divisor of \(F\) is zero,
\(F\) is an invertible regular function on a model neighbourhood
of \(C_0\setminus B\), so \(|F|=1\) on its tube.
In particular \(m_p=0\) for \(p\notin B\).
Some points of \(B\) can also have \(m_p=0\): zeros and poles
may reduce to the same point. They must not be discarded.

\emph{The source cutoff and every boundary term.}
For each \(p\in B\), shrink the collar in
\eqref{eq:p213-residual-lift-collar} and choose
\(0<a_p<b_p'<\delta_p\). Let
\(\eta_p(s)\) be a smooth function with values in \([0,1]\),
equal to one for \(s\leq a_p\) and zero for \(s\geq b_p'\).
Extend it by zero at the centre \(w_p=0\). Define \(\eta\) on
\(C^{\mathrm{an}}\) by this function on each tube in \(B\), and
by one elsewhere. It is globally CLD smooth. To see the gluing
explicitly, \(1-\eta_p\) is supported in the compact closed
disc \(\{|w_p|\leq e^{-a_p}\}\) contained in the ordinary open
tube. Its extension by zero outside that tube is smooth, and
there are only finitely many such functions.

The function \(\eta\) vanishes on neighbourhoods of all zeros
and poles of \(F\). Its support is compact because \(C\) is
proper. In particular, it is compactly supported in the
ordinary open set where \(F\) is an analytic unit.

Choose \(\varphi\in C_c^\infty((0,\ell))\) with
\(\int\varphi(u)\,\mathrm du=1\).
Choose \(\chi\in C_c^\infty(\mathbf R)\), equal to one near zero,
with support in \((-\varepsilon,\varepsilon)\), where
\[
 0<\varepsilon<
   \min_{\substack{p\in B\\m_p\neq0}} |m_p|a_p .
\]
If there are no nonzero \(m_p\), the desired principal relation
is already trivial. On the unit locus of \(F\), consider the
\((1,2)\)-form
\[
 \beta=\varphi(u)\chi(t)\,
        \mathrm d'u\wedge\mathrm d''u\wedge\mathrm d''t.
\]
The form \(\eta\beta\), extended by zero, is a smooth compactly
supported test form on \(X\): the \(u\)-support is contained
in a compact closed subannulus of \(A\), the curve factor is
proper, and \(\eta\) is zero near the removed divisor.

There is an exact equality of smooth forms
\begin{equation}
\label{eq:p213-residual-no-cutoff-error}
 \mathrm d'\eta\wedge\beta=0 .
\end{equation}
Indeed, the support of \(\mathrm d'\eta\) is contained in the
finitely many closed transition collars \(a_p\leq s_p\leq b_p'\).
If \(m_p\neq0\), then \(|t|>\varepsilon\) there and \(\chi(t)=0\).
If \(m_p=0\), \eqref{eq:p213-residual-lift-collar} says that
\(t\) is constant on an ordinary neighbourhood of every point
of the transition collar, so \(\mathrm d''t=0\).
Both arguments concern forms on the source, before pairing
with \(T\). No mass-nullity or current slicing is needed.

A direct differentiation gives
\[
 \mathrm d'\beta
  =\varphi(u)\chi'(t)\,
    \mathrm d'u\wedge\mathrm d''u\wedge
    \mathrm d't\wedge\mathrm d''t .
\]
The derivative of \(\varphi\) contributes zero.
On the support of \(\eta\mathrm d'\beta\), only the collars
with \(m_p\neq0\) can occur, and on these \(\eta=1\).
More explicitly, outside the tubes in \(B\) one has \(t=0\);
inside a tube, wherever \(\eta\neq0\), the exact formula
\(t=m_ps_p\) holds. Because \(\chi'\) vanishes near zero and
outside \((-\varepsilon,\varepsilon)\), each remaining support
is compactly contained in its individual collar \(V_p\).
Thus \(\eta\mathrm d'\beta\) is the finite sum of the zero
extensions of its restrictions to these disjoint source
collars. This is the step retaining individual \(p\)-weights
rather than only a sum on the tropical target.

Closedness and \eqref{eq:p213-residual-no-cutoff-error} give
\[
 0=\langle T,\mathrm d'(\eta\beta)\rangle
   =\sum_{\substack{p\in B\\m_p\neq0}}
       b_T(p)\int_0^\ell\varphi(u)\,\mathrm du
        \int_0^{\delta_p}
           m_p^2\chi'(m_ps)\,\mathrm ds .
\]
For either sign of \(m_p\),
\[
 \int_0^{\delta_p}m_p^2\chi'(m_ps)\,\mathrm ds
   =m_p\bigl(\chi(m_p\delta_p)-\chi(0)\bigr)=-m_p.
\]
This explicitly records the two differential factors and the
one change-of-variable factor; replacing \(m_p\) by \(|m_p|\)
in the final relation would be wrong. We have proved
\eqref{eq:p213-residual-principal-relation} for every \(f\).

\emph{The divisor functional and the canonical factors.}
The divisor group is the free abelian group on closed points,
so the finite numbers \(b_T(p)\) define a finite-valued
homomorphism, nonnegative on effective divisors. \emph{Finite}
here means finite on each divisor, which has finite support;
it does not assert a finite total measure on the set of all
closed points.

For completeness, the elementary proper-curve argument now
gives proportionality to degree. Choose an effective divisor
\(H\) with \(\deg H+\chi(C_0,\mathcal O_{C_0})>0\).
For every degree-zero divisor \(D\) and every integer \(n\),
Riemann--Roch gives a nonzero section of \(\mathcal O(H+nD)\);
hence \(H+nD\) is linearly equivalent to an effective divisor.
Positivity and principal-divisor annihilation imply
\(0\leq\lambda_T(H)+n\lambda_T(D)\) for both signs of \(n\).
Thus \(\lambda_T(D)=0\). Applying this to
\((\deg q)p-(\deg p)q\), for one fixed closed point \(q\),
gives \(b_T(p)=c\deg p\), with \(c\geq0\).
The analytic construction preceding this paragraph is the new
interface; this algebraic paragraph alone is not that interface.

As a \(K_p\)-analytic product polyannulus, \(V_p\) has standard
Gauss-plane calibration one in the primitive coordinates
\(u,s_p\). Viewed as a \(K\)-analytic space, its calibration
is multiplied by \([K_p:K]=d_p\), exactly
\cite[Remark 14.4.3, p.~410]{CLD25}; the finite-map formula is
\cite[Proposition 14.4.2, pp.~409--410]{CLD25}.
Therefore the canonical integration current has
\(b_{\delta_X}(p)=d_p\).

There is one source Gauss point over each \((u,s_p)\).
For example, irreducibility of the separable residue polynomial
defining \(\kappa(p)/k\) persists over the purely transcendental
residue field of a Gauss point; its unramified extension is a
field, not a product of \(d_p\) fields. Equivalently, the standard
\(K_p\)-torus skeleton maps bijectively to the \(K\)-torus
skeleton. This is also the finite-extension skeleton statement
in \cite[Remark 14.4.3]{CLD25}. The primitive lattice is unchanged
because the constant-field extension is unramified.

If the separate source-calibration representation is assumed,
the pushforward coefficient just computed consequently reads
off the coefficient on this single source Gauss plane. On each
of its top cells, injectivity of \(q_p\) and
\eqref{eq:p213-residual-branch-coefficient} identify that
coefficient with \(b_T(p)\). The canonical coefficient is \(d_p\),
so their ratio is \(b_T(p)/d_p=c\).
This uses an actual injective source skeleton in these explicit
coordinates, not compatibility only at generic full-rank points
of an unspecified refinement. The proof makes no inference
about any other cells or about \(T=c\delta_X\) on all of \(X\).
\end{proof}
